\documentclass[12pt,a4paper]{article}
\usepackage[utf8]{inputenc}
\usepackage[T1]{fontenc}
\usepackage[brazilian]{babel}
\usepackage{lmodern}
\usepackage[top=3cm,bottom=2cm,left=3cm,right=2cm]{geometry}
\usepackage{setspace}
\usepackage{hyperref}
\usepackage{graphicx}
\usepackage{enumitem}
\hypersetup{colorlinks=true,linkcolor=black,urlcolor=blue}
\setlength{\parindent}{1.25cm}
\setlength{\parskip}{6pt}
\sloppy
\onehalfspacing
\begin{document}
\begin{center}
{\large\textbf{Faculdade Descomplica}}\\[4pt]
{\large Ciência da Computação -- Projeto de Extensão V (PEX V)}\\[4pt]
Implementação e Intervenção Tecnológica\\[12pt]
{\Large\textbf{Relatório do Projeto de Extensão V}}\\[6pt]
{\large Vitrine Digital e Canal Direto de Atendimento para Microempreendimento de Marcenaria Artesanal}\\[12pt]
\begin{tabular}{rl}
Discente: & José Inamar de Medeiros Júnior \\
Período de Execução: & 2026 \\
Carga Horária Dedicada: & 40 horas \\
Repositório: & \url{https://github.com/EmiyaKiritsugu3/PEX_V} \\
Website: & \url{https://emiyakiritsugu3.github.io/PEX_V/} \\
\end{tabular}
\end{center}
\vspace{12pt}





\section{Sumário Executivo}

O presente Projeto de Extensão V teve como propósito intervir nos gargalos de captação, apresentação e comunicação de um microempreendimento artesanal do ramo de marcenaria de móveis planejados, operado informalmente em espaço residencial (garagem).

Diferente do cenário do Projeto de Extensão IV (focado na automação de processos administrativos e prescrições de treinos em uma academia comercial com Next.js, Prisma e IA), a presente demanda exigia uma abordagem de \textbf{Engenharia Frugal} e \textbf{Sustentabilidade Operacional (TCO Zero)}. O parceiro comunitário não possuía orçamento para infraestrutura de servidores, nem conhecimento técnico para gerenciar bancos de dados ou pipelines de compilação complexos.

Como intervenção tecnológica, desenvolveu-se uma plataforma web estática de página única (\textit{Single-Page Showcase}), focada em alta performance, carregamento instantâneo em redes móveis restritas (3G/4G) e conversão imediata para atendimento via WhatsApp. A solução eliminou qualquer dependência de backend pago, custos recorrentes de hospedagem e complexidade de manutenção, permitindo que o próprio artesão atualize seu portfólio substituindo arquivos locais de imagem.



\section{Contextualização e Diagnóstico da Situação-Problema}

A marcenaria em análise opera no modelo sob encomenda, focada na fabricação de armários, cozinhas compactas, bancadas de trabalho e nichos planejados. Embora o artesão detenha excelente domínio prático do ofício, o empreendimento sofria de severas restrições operacionais e mercadológicas:

\begin{enumerate}[leftmargin=*]
\item[\textbf{Dependência Exclusiva do Boca a Boca:}] A clientela restringia-se ao círculo imediato de vizinhança e indicações familiares.
\item[\textbf{Ausência de Referencial Visual Padronizado:}] Clientes solicitavam orçamentos sem conhecer o padrão de acabamento, exigindo visitas presenciais demoradas ou envio desordenado de fotos dispersas na galeria do celular.
\item[\textbf{Assimetria de Informações na Triagem:}] As conversas iniciais demandavam tempo excessivo do marceneiro para explicar modelos básicos, medidas aproximadas e prazos, subtraindo horas produtivas da bancada de trabalho.
\item[\textbf{Vulnerabilidade Financeira:}] Impossibilidade de arcar com mensalidades de plataformas SaaS (Shopify, Wix, construtores pagos) ou custos de servidores dedicados em nuvem.
\end{enumerate}



\section{Análise Crítica dos Processos}

\subsection{Processo Operacional Anterior}
\begin{enumerate}[leftmargin=*]
\item O cliente entrava em contato por telefone ou mensagem sem contexto prévio;
\item O marceneiro interrompia a fabricação manual para responder e buscar fotos antigas no celular;
\item Negociações eram conduzidas sem uma lista prévia de modelos fabricáveis;
\item Dificuldade de estabelecer credibilidade profissional inicial com novos clientes fora do bairro.
\end{enumerate}

\subsection{Principais Problemas Identificados}
\begin{itemize}[leftmargin=*]
\item[\textbf{Custo de Oportunidade Elevado:}] O tempo gasto em triagens repetitivas reduzia o tempo de fabricação e montagem;
\item[\textbf{Falta de Vitrine Centralizada 24/7:}] Perda de oportunidades comerciais fora do horário de trabalho;
\item[\textbf{Fragmentação Visual:}] Comunicação amadora e sem identidade consolidada.
\end{itemize}



\section{Análise de Causa Raiz -- Técnica dos 5 Porquês}

\begin{itemize}[leftmargin=*]
\item[\textbf{Problema Central:}] Baixo volume de orçamentos qualificados e sobrecarga de atendimento manual.
\end{itemize}
\begin{enumerate}[leftmargin=*]
\item \textit{Por que os orçamentos demoravam a fechar?} O cliente não tinha clareza prévia sobre modelos, medidas e acabamentos oferecidos.
\item \textit{Por que o cliente não tinha acesso a essas informações?} Não existia um catálogo ou vitrine digital acessível publicamente.
\item \textit{Por que o artesão não utilizava uma vitrine digital?} As ferramentas convencionais do mercado impõem barreiras financeiras e técnicas de manutenção.
\item \textit{Por que essas ferramentas são inacessíveis para a realidade da garagem?} Exigem custos recorrentes (mensalidades/servidores) e conhecimentos de TI para atualização que o marceneiro não possui.
\item \textit{Por que não se adotou uma solução simplificada antes?} Presunção de que soluções web exigem bancos de dados complexos e servidores contínuos.
\end{enumerate}
\begin{itemize}[leftmargin=*]
\item[\textbf{Causa Raiz Identificada:}] Falta de uma solução tecnológica sob medida, regida pela frugalidade, que conecte a vitrine ao canal direto de conversão (WhatsApp) sem custos e sem atrito operacional.
\end{itemize}



\section{Fatores-Chave para a Intervenção Tecnológica}

\subsection{Fator Humano e Cognitivo}
A equipe operacional consiste em uma única pessoa, responsável tanto pela produção física dos móveis quanto pela administração. A interface precisa ser autoexplicativa, sem exigir painel administrativo (CMS) complexo ou logins.

\subsection{Fator Financeiro (TCO Zero)}
O projeto adota a premissa de custo contínuo igual a zero. Utiliza arquitetura estática publicada em plataforma de borda gratuita (GitHub Pages), sem necessidade de banco de dados relacional pago ou serviços em segundo plano.

\subsection{Fator de Mobilidade e Conectividade}
A grande maioria dos clientes acessa a plataforma via smartphone em conexões de baixa largura de banda. A página completa pesa $\sim$42 KB (excluindo imagens), sem bloqueios de renderização por bibliotecas pesadas.

\subsection{Fator de Sustentabilidade Técnica (Green Software Engineering)}
Evitou-se o desperdício computacional típico de SPAs pesadas (Node runtime, SSR, ORM). O código HTML/CSS puro roda diretamente no cliente sem consumo ocioso de CPU em nuvem.



\section{Proposição e Justificativa Arquitetural}

Em contraste com o PEX IV (onde a complexidade pedagógica e biométrica exigia persistência relacional com Prisma e sugestões via IA), a engenharia aplicada no PEX V baseou-se nos princípios \textbf{YAGNI} (\textit{You Aren't Gonna Need It}) e \textbf{KISS} (\textit{Keep It Simple, Stupid}), oriundos da programação extrema \cite{beck1999}.

\subsection{Decisão Arquitetural: Arquitetura Estática Desacoplada}
\begin{itemize}[leftmargin=*]
\item[\textbf{Arquivo Único (\texttt{index.html}):}] CSS e JS estritamente inline, garantindo carregamento atômico em uma única requisição HTTP;
\item[\textbf{Galeria Tolerante a Falhas:}] Estrutura concebida para exibir placeholders informativos caso as fotos estejam ausentes, e substituição gradual por fotos locais na pasta \texttt{fotos/} com tratativa nativa de erro (\texttt{onerror});
\item[\textbf{Encaminhamento Semântico via WhatsApp (\texttt{wa.me}):}] Botões contextuais para cada modelo com mensagens pré-codificadas (\texttt{encodeURIComponent}), entregando o cliente já qualificado para o fechamento da venda;
\item[\textbf{Validação Automatizada Local (\texttt{check.py}):}] Script em Python puro para verificação estática de integridade (links, tags, integridade de caminhos de imagem e ausência de bloqueios em ambiente local \texttt{file://}).
\end{itemize}



\section{Estrutura do Sistema e Implementação}

\begin{verbatim}
projeto_de_extensão_V/
+-- index.html                     # Vitrine completa (HTML5 + CSS responsivo + JS inline)
+-- check.py                       # Verificador automatizado de integridade
+-- .gitignore                     # Governança de repositório
+-- fotos/                         # Registro fotográfico real da oficina
|   +-- oficina-bancada.jpg
|   +-- oficina-serra.jpg
|   \-- oficina-paineis.jpg
\-- docs/
    +-- relatorio-extensao-V.md    # Documentação acadêmica formal (este relatório)
    \-- superpowers/               # Especificação técnica e plano de engenharia
        +-- specs/2026-09-16-galeria-moveis-planejados-design.md
        +-- specs/2026-09-21-pex-v-nota-academica-design.md
        \-- plans/2026-09-16-galeria-moveis-planejados.md
\end{verbatim}

\subsection{Principais Seções Implementadas}
\begin{enumerate}[leftmargin=*]
\item[\textbf{Hero Section:}] Proposta de valor direta e chamada principal de ação;
\item[\textbf{Sobre:}] Apresentação da procedência artesanal e atendimento direto;
\item[\textbf{Galeria de Modelos:}] 6 nichos padronizados com mensagens específicas de orçamento;
\item[\textbf{Fluxo de 4 Passos:}] Explicação didática da jornada do cliente (Conversa $\to$ Cotação $\to$ Produção $\to$ Entrega/Instalação);
\item[\textbf{Canal de Contato Flutuante e Fixo:}] Garantia de acessibilidade permanente ao botão de orçamento.
\end{enumerate}



\section{Metodologia e Cronograma de Execução}

O desenvolvimento do projeto foi dividido em etapas organizadas conforme a carga horária total de 40 horas dedicadas ao componente curricular, seguindo a sequência natural de uma intervenção extensionista: compreender a realidade, projetar a solução, construir, validar e implantar com documentação.

\subsection{Levantamento e Diagnóstico (8 horas)}

\paragraph{Atividades Realizadas}
\begin{itemize}[leftmargin=*]
\item Visitas à marcenaria e registro fotográfico do espaço de trabalho (oficina em garagem);
\item Conversas com o artesão responsável para mapear a rotina de atendimento e os gargalos de captação de clientes;
\item Levantamento das restrições financeiras (ausência de orçamento para SaaS/servidores) e técnicas (sem conhecimento de TI para manutenção);
\item Formalização da parceria por meio do Termo de Autorização para realização das atividades extensionistas.
\end{itemize}

\subsection{Design e Prototipação (8 horas)}

\paragraph{Atividades Realizadas}
\begin{itemize}[leftmargin=*]
\item Elaboração da especificação técnica (\path{docs/superpowers/specs/2026-09-21-pex-v-nota-academica-design.md});
\item Definição da identidade editorial artesanal (tipografia Fraunces + Outfit + DM Mono, paleta madeira/linho);
\item Prototipação da estrutura de página única: Hero, Sobre, Galeria de Modelos (6 nichos), Fluxo em 4 passos e Contato;
\item Desenho dos esquemas técnicos em SVG inline como placeholders tolerantes à ausência de fotos.
\end{itemize}

\subsection{Desenvolvimento da Vitrine (12 horas)}

\paragraph{Atividades Realizadas}
\begin{itemize}[leftmargin=*]
\item Implementação do \texttt{index.html} em arquivo único (HTML5 + CSS responsivo + JS inline), sem dependências de build;
\item Implementação do encaminhamento semântico via WhatsApp (\texttt{wa.me} + \texttt{encodeURIComponent}), recurso oficial de conversa sem salvamento de contato \cite{whatsapp}, com mensagens pré-codificadas por modelo;
\item Tratamento defensivo para número não configurado (ocultação dos botões + aviso \texttt{\#sem-contato});
\item Pasta \texttt{fotos/} para atualização do portfólio pelo próprio artesão via substituição de arquivos locais.
\end{itemize}

\subsection{Validação e Testes (6 horas)}

\paragraph{Atividades Realizadas}
\begin{itemize}[leftmargin=*]
\item Desenvolvimento do verificador automatizado \texttt{check.py} (constante WhatsApp única, ausência de \texttt{fetch}, \texttt{data-msg} em todos os \texttt{wa-link}, integridade de \texttt{fotos/});
\item Validação de abertura direta em \texttt{file://} (sem servidor), garantindo funcionamento offline local;
\item Testes de responsividade com ênfase em telas de 360px (micro-ritmo) e navegação mobile;
\item Inserção das 3 fotos reais da oficina e revalidação integral.
\end{itemize}

\subsection{Implantação e Documentação (6 horas)}

\paragraph{Atividades Realizadas}
\begin{itemize}[leftmargin=*]
\item Publicação da vitrine no GitHub Pages \cite{githubpages} a partir da branch \texttt{main} (\url{https://emiyakiritsugu3.github.io/PEX_V/});
\item Versionamento do trabalho em repositório público com Pull Request revisado e mergeado;
\item Redação deste relatório e organização da documentação acadêmica em \texttt{docs/};
\item Preenchimento e conferência do Termo de Autorização junto à organização parceira.
\end{itemize}



\section{Resultados Obtidos}

A intervenção entregou uma vitrine digital funcional, publicada e validada, que centraliza a apresentação dos modelos fabricáveis e qualifica o primeiro contato via WhatsApp. Mesmo em caráter de demonstração acadêmica, foi possível estruturar:

\begin{itemize}[leftmargin=*]
\item Vitrine de página única publicada em URL pública, com carregamento instantâneo;
\item Galeria de 6 modelos padronizados (guarda-roupa, cozinha, rack/painel, home office, gabinete de banheiro, estante) com mensagens de orçamento específicas;
\item Registro fotográfico real da oficina (bancada, serra, painéis) integrado à seção Sobre (Figura \ref{fig:oficina});
\item Canal de conversão direto via WhatsApp, com fallback informativo quando o número comercial ainda não está configurado;
\item Custo recorrente zero e manutenção delegável ao próprio artesão.
\end{itemize}

\subsection{Resultados Esperados}
\begin{itemize}[leftmargin=*]
\item Redução de 70\% no tempo gasto explicando dimensões e modelos básicos durante a triagem inicial;
\item Aumento da taxa de conversão através de mensagens pré-formatadas no WhatsApp;
\item Estabelecimento de uma identidade profissional para o empreendimento perante o mercado local;
\item Zero custo recorrente para o artesão ao longo do ciclo de vida da aplicação.
\end{itemize}

\subsection{Próximos Passos}
\begin{enumerate}[leftmargin=*]
\item Coleta e inserção de fotos reais de projetos finalizados (móveis entregues) na pasta \texttt{fotos/};
\item Substituição da constante \texttt{NUMERO\_WHATSAPP} pelo número comercial definitivo;
\item Coleta de dados quantitativos de novos clientes contactados após 30 dias de implantação.
\end{enumerate}

\begin{figure}[h]
\centering
\includegraphics[width=0.32\textwidth]{fotos/oficina-bancada.jpg}\hfill
\includegraphics[width=0.32\textwidth]{fotos/oficina-serra.jpg}\hfill
\includegraphics[width=0.32\textwidth]{fotos/oficina-paineis.jpg}
\caption{Registro fotográfico da oficina parceira: bancada com ferramental, área de corte com serra e chapas em produção, e estoque de painéis.}
\label{fig:oficina}
\end{figure}



\section{Alinhamento com a Área de Ciência da Computação}

O projeto permitiu aplicar conhecimentos fundamentais do curso de Ciência da Computação, incluindo:

\paragraph{Desenvolvimento Web}
Construção de página responsiva com HTML5 semântico, CSS moderno (grid, variáveis, media queries) e JavaScript client-side para composição dinâmica de links, conforme a documentação de referência das tecnologias web \cite{mdn}.

\paragraph{Engenharia de Software}
Aplicação dos princípios YAGNI e KISS, validação automatizada estática (\texttt{check.py}), versionamento com Git/GitHub via Pull Request e publicação contínua pelo GitHub Pages.

\paragraph{Interação Humano-Computador (IHC)}
Projeto mobile-first para público de baixa conectividade, hierarquia visual editorial, tratamento de estados vazios (placeholders SVG, aviso de contato indisponível) e jornada de conversão em 4 passos, em consonância com os princípios de engenharia de usabilidade \cite{nielsen1993}.

\paragraph{Redes e Infraestrutura}
Hospedagem estática em borda (entrega HTTP, cache e HTTPS delegados à plataforma), página completa em $\sim$42 KB para operação em redes 3G/4G restritas.

\paragraph{Sustentabilidade Computacional (Green Software)}
Arquitetura sem backend ocioso: zero consumo de CPU em nuvem fora do momento do acesso, em consonância com os princípios de engenharia de software sustentável \cite{gsf}, em contraste com SPAs pesadas com runtime e SSR contínuos.



\section{Conclusões e Aprendizados}

A realização deste Projeto de Extensão V proporcionou experiência prática na análise de um microempreendimento real e no desenvolvimento de uma solução tecnológica sob medida para suas restrições --- financeiras, técnicas e operacionais. O contato direto com a marcenaria permitiu compreender que, para pequenos negócios informais, a melhor engenharia é a que o próprio artesão consegue sustentar.

O projeto demonstrou, em caráter de demonstração acadêmica, como uma intervenção frugal pode profissionalizar a apresentação comercial e organizar o atendimento sem impor custos ou complexidade ao parceiro comunitário.

\paragraph{Principais Aprendizados}
\begin{itemize}[leftmargin=*]
\item Desenvolvimento web estático performático e responsivo;
\item Engenharia frugal e análise de custo total (TCO) como critério arquitetural;
\item Validação automatizada e testes de compatibilidade local (\texttt{file://});
\item Publicação e versionamento com GitHub (PR, merge, Pages);
\item Levantamento de requisitos com cliente não técnico e formalização via Termo de Autorização;
\item Escrita técnica acadêmica com evidências comprobatórias.
\end{itemize}

Além do crescimento técnico, o projeto reforçou a importância da extensão universitária como ponte entre o conhecimento acadêmico e as necessidades concretas da comunidade local, conforme preconizam as diretrizes curriculares da extensão \cite{resolucao2018}.



\section{Evidências Comprobatórias}

\paragraph{Repositório do Código-Fonte}
https://github.com/EmiyaKiritsugu3/PEX_V

\paragraph{Website do Projeto (GitHub Pages)}
https://emiyakiritsugu3.github.io/PEX_V/

\paragraph{Estrutura do Projeto}
\begin{itemize}[leftmargin=*]
\item \texttt{index.html} $\to$ vitrine completa (HTML5 + CSS responsivo + JS inline);
\item \texttt{check.py} $\to$ verificador automatizado de integridade;
\item \texttt{fotos/} $\to$ \texttt{oficina-bancada.jpg}, \texttt{oficina-serra.jpg}, \texttt{oficina-paineis.jpg} (registro real da oficina);
\item \texttt{docs/relatorio-extensao-V.md} $\to$ documentação acadêmica formal (este relatório);
\item \texttt{docs/superpowers/} $\to$ especificação técnica e plano de engenharia;
\item \texttt{Termo de Autorização} $\to$ preenchido e conferido com a organização parceira (JS Planejados -- Móveis sob medida, São João do Sabugi/RN); mantido fora do repositório público por conter dados sensíveis (CNPJ), arquivado junto à documentação da disciplina.
\end{itemize}

\paragraph{Interfaces Desenvolvidas}
\begin{itemize}[leftmargin=*]
\item Hero com proposta de valor e chamada para WhatsApp;
\item Seção Sobre com procedência artesanal, esquema técnico e fotos da oficina;
\item Galeria de 6 modelos sob medida com orçamento contextual por modelo;
\item Fluxo de 4 passos (Conversa $\to$ Cotação $\to$ Produção $\to$ Entrega/Instalação);
\item Contato fixo e botão flutuante de WhatsApp com fallback informativo.
\end{itemize}

A apresentação deste relatório segue a ABNT NBR 14724 \cite{abnt14724}, e as referências foram elaboradas conforme a ABNT NBR 6023 \cite{abnt6023}.

\begin{thebibliography}{9}
\bibitem{abnt14724} ASSOCIAÇÃO BRASILEIRA DE NORMAS TÉCNICAS. \textbf{NBR 14724}: informação e documentação --- trabalhos acadêmicos --- apresentação. Rio de Janeiro, 2024.

\bibitem{abnt6023} ASSOCIAÇÃO BRASILEIRA DE NORMAS TÉCNICAS. \textbf{NBR 6023}: informação e documentação --- referências --- elaboração. Rio de Janeiro, 2018.

\bibitem{resolucao2018} BRASIL. Ministério da Educação. Conselho Nacional de Educação. Resolução CNE/CES nº 7, de 18 de dezembro de 2018. Estabelece as Diretrizes para a Extensão na Educação Superior Brasileira. Diário Oficial da União, Brasília, 19 dez. 2018.

\bibitem{beck1999} BECK, K. \textbf{Extreme Programming Explained}: embrace change. Reading: Addison-Wesley, 1999.

\bibitem{nielsen1993} NIELSEN, J. \textbf{Usability Engineering}. San Francisco: Morgan Kaufmann, 1993.

\bibitem{gsf} GREEN SOFTWARE FOUNDATION. \textbf{Principles of Green Software Engineering}. Disponível em: \url{https://principles.green/}. Acesso em: 26 set. 2026.

\bibitem{githubpages} GITHUB. \textbf{GitHub Pages Documentation}. Disponível em: \url{https://docs.github.com/pages}. Acesso em: 26 set. 2026.

\bibitem{mdn} MOZILLA. \textbf{MDN Web Docs}. Disponível em: \url{https://developer.mozilla.org/}. Acesso em: 26 set. 2026.

\bibitem{whatsapp} WHATSAPP. \textbf{Click to Chat}: converse com uma pessoa sem ter o número salvo. Disponível em: \url{https://faq.whatsapp.com/}. Acesso em: 26 set. 2026.
\end{thebibliography}



\end{document}
